% ---
% id: transmitter/rem:boundary
% kind: rem
% title: "Scope of a verified certificate"
% summary: >-
%   A verified certificate buys the interchange and the source identification
%   of [[thm:certificate-sufficiency]] outright, and in block form the reached-state statement of
%   [[lem:matched-pinching]]. It is no security reduction, restricts no
%   consumer's feasible set and decides nothing about retention.
% prominence: standard
% status: accepted
% uses: [thm:certificate-sufficiency, def:srcert, lem:matched-pinching, def:source-segment]
% mentions: [ass:src, ass:seq, rem:xreg, def:source-spec, req:attestation-record, def:source-link]
% assumes: []
% refs: []
% provenance: newly-derived
% lean:
% ---

\begin{remark}
\label{rem:boundary}
Theorem~\ref{thm:certificate-sufficiency}\,(iii) is a statement about the \emph{source}: it
says that Alice's measurement commutes with every attack of the admitted
form --- every channel on the \tref[virtual transmitted system]{transmitted system} and the adversary's
registers, which by Assumption~\ref{ass:seq} is every operation of the
round up to the announcements (Remark~\ref{rem:xreg}) --- and that the
declared trusted-source data describe the virtual preparation (the
specified one exactly when the \tref{source segment} is empty) against every such
attack. It is not a security reduction, and it does not say that the
consuming engine's optimization is sound; that argument --- granting the
adversary a purification of Alice's retained system and the
\tref[virtual transmitted system]{transmitted system}, restricting to the declared feasible set, accumulating over rounds
--- lives in the engine's claims, which import Assumption~\ref{ass:src} for
it. What this module owes depends on the form of the certificate
(Theorem~\ref{thm:certificate-sufficiency}\,(i)). In pure and in block form it is the
decided conditions of the declared certificate, those on $\Psi$ of its form
together with (SR1) and (SR2), and the attestation (SR3). The conditions on $\Psi$
are Hermiticity, $\Psi^2=\Psi$ and $\tr\Psi=1$ in pure form, and Hermiticity
and positive semidefiniteness of every block and $\sum_t\tr\Psi_t=1$ in block
form. For an ensemble presentation it is exactly (E1) to (E3), decided on
the declared blocks, and the attestation (SR3); (SR1) and (SR2) then hold of
the certificate the presentation defines as theorems and are decided by
nothing.

\emph{What a verified certificate buys.} A verified certificate, one whose
decided conditions (Theorem~\ref{thm:certificate-sufficiency}\,(i)) hold, buys
\eqref{eq:srmain} and \eqref{eq:srident} outright: the commutation, and the
identification of the declared trusted-source data with the virtual source,
as identities for every channel of the admitted form, with no residue and no
qualification. That is what the exact reading of (SR1)--(SR2) is for, and it
is what a consumer of the licence record, the record of the verified
certificate that the artefact carries (Requirement~\ref{req:attestation-record}),
is entitled to read the record as asserting. What the certificate does
\emph{not} buy is (SR3), which no function of the declared data decides
(Assumption~\ref{ass:src}); it is (SR3), with Assumption~\ref{ass:seq}, that
makes the attack of a round a channel of the admitted form.

\emph{What the block form adds.} A certificate declared in block form
(Definition~\ref{def:srcert}) buys one thing more, and it is a statement about
\emph{reached} states: Lemma~\ref{lem:matched-pinching}, that every state an
admitted attack reaches from that certificate is block diagonal along the
declared sectors with $A$-marginal $\tau_A$, and that pinching $A$ after the
attack is the same operation as pinching the virtual signal before it. That is
what a consumer may read the record as asserting of the declared \tref{sector
structure}, and it is all of it. \emph{It does not buy a restriction of any
consumer's feasible set.} A feasible set is a set of operators constrained to
have the marginal $\tau_A$, not a set of states reached from a certificate,
and a positive operator with that marginal need not be block diagonal
(Lemma~\ref{lem:matched-pinching}\,(iii)). A consumer that minimises over
block-diagonal operators alone does so under a theorem of its own module, with
its own preconditions on declared data, and neither this module's licence
record nor the presentation of $\tau_A$ and $\{M_x\}$ as block lists is that
theorem (Definition~\ref{def:source-spec}, ``What the sector structure is
not'').

\emph{Links supplied by a theorem.} The same boundary, in the other
direction, governs a \tref{source link} whose channel is supplied by a named claim
of a submodule rather than as data (a \tref{theorem constructor},
Definition~\ref{def:source-link}\,(c2)). For such a link, what the artefact
carries about its source identity \eqref{eq:srcmap} is the link's \tref{obligation
record}, which travels with the link in the source segment
(Definition~\ref{def:source-segment}\,(S4)) and not in the licence record: a
named claim and the assumptions the obligation carries, never a verdict this
module decided, and there is nothing of that identity a consumer could
re-decide.
\end{remark}
